\chapter*{Summary}
\label{ch:summary}
\addcontentsline{toc}{chapter}{\nameref{ch:summary}}

\vspace*{-1cm}
\epigraph{"Everything is a Lagrangian submanifold"}{Alan Weinstein}

The main goal of this thesis is to classify Lagrangian fibrations over compact
surfaces. The base space of such a fibration comes equipped with an
\emph{integral affine structure} which is an invariant of the Lagrangian
fibration. We classify the Lagrangian fibrations over a compact surface $B$ by
first classifying all the possible integral affine structures on $B$ and then
classifying all of the Lagrangian fibrations corresponding to a fixed integral
affine structure. Remarkably, if $B$ is simply connected the Lagrangian
fibrations over $B$ are determined entirely by this integral affine structure.

In Chapter \ref{ch:prelim}, we develop the basic notions of affine geometry
which belongs to a more general class of geometries called
$(G,X)$\emph{-structures} which were studied by Thurston, Ehresmann and many
others \cite{thurston1980geometry}, \cite{Ehresmann1936}. A common theme when
classifying the geometric structures on a manifold $M$, is that the fundamental
group plays an important role in which geometric structures $M$ can be equipped
with. To this end, we develop the notion of representation theory for affine
manifolds. We also describe how an affine manifold can be equipped with a
canonical flat, torsion-free connection and how geodesic completeness of this
connection relates to the affine structure. Finally, we introduce the notion of
an integral affine manifold and give some properties and examples of these
objects.

In Chapter \ref{ch:nilpotent-hol}, we describe how the condition of the
fundamental group being nilpotent relates to completeness of the affine
structure. A common theme in the study of affine manifolds is that the
fundamental group determines many properties of the affine structure. It is
conjectured by Auslander in \cite{AUSLANDER1964131} that the fundamental group
of a complete affine manifold should be solvable. This was shown to be true up
to dimension six in \cite{abels2020auslanderconjecturedimension7}. We
intertwine the ideas of this conjecture with another open problem called the
\emph{Markus conjecture}. This states that for affine manifolds, having
parallel volume is equivalent to being complete. The main result of this
chapter is that the Markus conjecture holds if the fundamental group is
nilpotent as first shown in \cite{nilpotent-holonomy}.

Chapter \ref{LagrangianFibrations} introduces Lagrangian fibrations and we
describe some of their properties and give some examples. We show how the base
space of a regular Lagrangian torus fibration inherits the structure of an
integral affine manifold. The fibration allows us to realise the integral
affine structure on the base space via so called action-angle coordinates. The
action coordinates will determine the integral affine structure on the base and
the angle coordinates will give coordinates on the fibres of the fibration. We
will look at the obstructing to our manifold admitting \emph{global}
action-angle coordinates which was studied in detail by Duistermaat in
\cite{ActionAngle}.

In Chapter \ref{ch:class-IAS}, we show that the only compact surfaces that can
admit an integral affine structure are the Klein bottle and the torus. Using
the notion of completeness for affine structures and exploiting that the
fundamental group of the torus is nilpotent, we can reduce the problem of
classifying the integral affine structures on these surfaces to a problem in
linear algebra and group theory. Classifying the injective homomorphisms from
the fundamental group of these surfaces to the integral affine transformation
group of $\R^2$ such that the quotient is a smooth manifold, is shown to be
equivalent to classifying the integral affine structures on these surfaces.

In Chapter \ref{ch:classify-fibr}, we classify the Lagrangian fibrations over
each of these integral affine surfaces. The integral affine structure
determines the local nature of such a Lagrangian fibration through the so
called \emph{symplectic reference fibration}. The symplectic reference
fibration can be glued together in various ways to produce all Lagrangian
fibrations over the base space that induce the given integral affine structure.
Exploiting completeness of the integral affine structures on these surfaces, we
can classify these Lagrangian fibrations topologically using techniques from
group cohomology and sheaf theory. We explicitly calculate the invariants for
these Lagrangian fibrations and show how we can construct them from their
invariants using Luttinger surgery.

